\documentclass[11pt,a4paper]{article}

\usepackage[T1]{fontenc}
\usepackage[utf8]{inputenc}
\usepackage[spanish,es-noshorthands]{babel}
\usepackage{lmodern}
\usepackage[margin=2.4cm]{geometry}
\usepackage{amsmath,amssymb}
\usepackage{booktabs}
\usepackage{longtable}
\usepackage{array}
\usepackage{multirow}
\usepackage{textcomp}
\providecommand{\degree}{\ensuremath{^\circ}}
\usepackage{xcolor}
\usepackage{graphicx}
\usepackage[colorlinks=true,linkcolor=azulunal,urlcolor=azulunal,citecolor=azulunal]{hyperref}
\usepackage{titlesec}
\usepackage{enumitem}

\definecolor{azulunal}{RGB}{20,70,130}
\definecolor{grisclaro}{RGB}{245,245,247}
\definecolor{rojoalerta}{RGB}{170,40,40}

\titleformat{\section}{\large\bfseries\color{azulunal}}{\thesection}{0.7em}{}
\titleformat{\subsection}{\normalsize\bfseries}{\thesubsection}{0.6em}{}

% ---------------------------------------------------------------------------
% Enlaces al repositorio. No se usan rutas absolutas en ningun sitio.
% ---------------------------------------------------------------------------
\newcommand{\ghrepo}{https://github.com/gustavop-dev/humerous-detection-fat-supress}
\newcommand{\ghbranch}{geometry-exercises}
\newcommand{\ghblob}{\ghrepo/blob/\ghbranch}
\newcommand{\gh}[1]{\href{\ghblob/#1}{\texttt{#1}}}
\newcommand{\ghl}[2]{\href{\ghblob/#1}{#2}}

\newcommand{\Rd}{\ensuremath{R/d}}
\newcommand{\Ot}{\ensuremath{O(3,1)}}

\title{\vspace{-1.5cm}\textbf{Representación geométrica del húmero y transporte\\
de las estructuras musculares mediante \Ot}\\[0.35cm]
\large Implementación de los ejercicios propuestos y análisis de resultados}
\author{Gustavo Adolfo Pérez\\
\small Universidad Nacional de Colombia --- Tópicos en Geometría Computacional\\
\small Ejercicios propuestos por el profesor Marco Paluszny}
\date{\today}

\begin{document}
\maketitle

\begin{abstract}
\noindent
Se implementan los ejercicios de geometría proyectiva propuestos sobre el pipeline de
segmentación del húmero en MRI axial con saturación de grasa. El objetivo final es
transportar las mallas de los músculos del manguito rotador desde un húmero de referencia
al húmero de un paciente, usando el par (esfera articular, eje diafisario) como única
correspondencia y una transformación del grupo \Ot{} que preserva esferas y manda rectas
en rectas.

El núcleo matemático se verifica contra el resultado publicado en el enunciado y pasa
16 de 16 pruebas. El transporte completo se ejecuta sobre 3 datasets, con verificación
numérica de que la transformación lleva la esfera del referente exactamente sobre la del
paciente ($5\times10^{-14}$ mm) y el eje sobre el eje ($5\times10^{-13}$ mm).

El hallazgo principal es negativo y es el más útil del trabajo: la aparente violación de
la condición $R_1/d_1 = R_2/d_2$ ---que justificaría la necesidad de \Ot{}--- resultó ser
en su mayor parte un artefacto de la cobertura axial limitada de los volúmenes de
resonancia, no una diferencia anatómica. Se cuantifica el sesgo, se corrige, y la
corrección se valida de forma cruzada: mejora la fidelidad del hueso transportado de
4.4 mm a 3.3 mm de RMS contra datos que no intervinieron en derivarla.
\end{abstract}

\vspace{-0.3cm}
\begin{center}
\colorbox{grisclaro}{\parbox{0.94\textwidth}{\small
\textbf{Repositorio:} \url{\ghrepo} \quad(rama \texttt{\ghbranch})\\[2pt]
\textbf{Este informe en PDF:} \ghl{informe/informe.pdf}{\texttt{informe/informe.pdf}}
\quad\textbf{fuente:} \ghl{informe/informe.tex}{\texttt{informe/informe.tex}}\\[2pt]
\textbf{Demostraciones de papel:} \ghl{excercises/DEMOSTRACIONES.md}{\texttt{excercises/DEMOSTRACIONES.md}}
}}
\end{center}

% ===========================================================================
\section{Contexto y objetivo}

El enunciado (30 diapositivas) converge en una sola actividad: extraer de un húmero de
referencia, dado como malla \texttt{.stl}, el par (esfera, eje); extraer el mismo par del
húmero de un paciente a partir de los arcos circulares detectados en su resonancia;
construir la transformación que lleva un par en el otro; y usarla para transportar las
mallas de los músculos del manguito rotador.

La esfera aproxima la superficie articular de la cabeza humeral y el eje aproxima el eje
del brazo. Anatómicamente el centro de la esfera es \emph{excéntrico} respecto al eje, y esa
excentricidad $d$ es la cantidad clave: la razón adimensional \Rd{} decide qué grupo de
transformaciones hace falta.

El interés va más allá del ejercicio. La presentación del proyecto para el congreso UEMS
(\ghl{presentation/GUION.md}{\texttt{presentation/GUION.md}}) asume una regla fija de
sectores angulares para asignar cada tendón ---subescapular $0$--$45\degree$, supraespinoso
$45$--$135\degree$, infraespinoso $135$--$225\degree$, redondo menor
$225$--$270\degree$--- y el propio guion la reconoce como hipótesis sin validar.
Transportar las mallas musculares reales al húmero del paciente es justamente lo que
convertiría esa regla en una cantidad medida por paciente.

\subsection{El nudo técnico}

Una similaridad $T(x) = (sR)x + b$ multiplica todas las distancias por el mismo factor $s$.
Manda esferas en esferas y rectas en rectas, pero \emph{preserva toda cantidad adimensional
construida con distancias}, en particular \Rd. Por tanto solo puede llevar un par
(esfera, eje) exactamente sobre otro si

\[
\frac{R_1}{d_1} \;=\; \frac{R_2}{d_2}.
\]

Cuando la igualdad falla hay que salir del grupo de similaridades. El grupo \Ot{} ---las
matrices $M$ con $M^{\mathsf T} \eta M = \eta$, $\eta = \mathrm{diag}(1,1,1,-1)$--- preserva la
cuádrica $x^2+y^2+z^2-w^2=0$ de $P^3$ (esto es, la esfera) y manda rectas en rectas, pero
\emph{no} preserva el centro de la esfera, así que sí puede alterar \Rd.

% ===========================================================================
\section{Ejercicios y dónde está cada uno}

\begin{longtable}{@{}p{1.6cm} p{6.4cm} p{7.6cm}@{}}
\toprule
\textbf{Diap.} & \textbf{Ejercicio} & \textbf{Dónde está resuelto} \\
\midrule
\endhead
4 & $R \in O(3) \Rightarrow \mathrm{diag}(R,1) \in \Ot$ &
Demostración en \gh{excercises/DEMOSTRACIONES.md}; \texttt{embed\_rotation} \\
\addlinespace
11--12 & Construir $T$ de dos pares esfera/eje; verificar $R$ ortogonal &
\texttt{similarity\_from\_pairs} en \gh{Geometry/similarity.py} \\
\addlinespace
13 & Ilustrar $R_1/d_1 = R_2/d_2$ cumplida y violada &
Demostración y \texttt{pair\_similarity\_invariant} \\
\addlinespace
17a & Siempre existe $T$ con $S_1 \to S_2$ y $L_1 \parallel L_2$ &
Demostración constructiva; \texttt{similarity\_axis\_parallel} \\
\addlinespace
17b & Por qué no se logra $T(L_1) = L_2$ &
Demostración (argumento del cilindro) \\
\addlinespace
18 & Verificar que $M$ preserva la esfera y manda $A_1,B_1 \to A_2,B_2$ &
\texttt{test\_boost\_ground\_truth} en \gh{Geometry/selftest.py} \\
\addlinespace
19 & Script en Python que construye $M \in \Ot$ &
\texttt{boost\_fixing\_sphere} en \gh{Geometry/lorentz.py} \\
\addlinespace
24 & Contraejemplo (vale doble) &
\texttt{build\_non\_similar\_example} y demostración \\
\addlinespace
25a & Detectar arcos circulares: secuencia estructurada &
Ya implementado en \gh{humerus\_boundary\_analysis.py}; documentado en 9 pasos \\
\addlinespace
25b & Estrategia para detectar el eje del húmero &
\gh{Geometry/sphere\_axis.py} \\
\addlinespace
29 & Actividad completa: transporte de los músculos &
\gh{Geometry/transfer.py} \\
\bottomrule
\end{longtable}

% ===========================================================================
\section{Software implementado}

Se añadió el paquete \ghl{Geometry}{\texttt{Geometry/}} y el ejecutable
\gh{exercises\_geometry.py}. \textbf{No se introdujo ninguna dependencia nueva} y
\textbf{no se modificó} el pipeline de segmentación existente: todo el trabajo nuevo
consume los CSV que ese pipeline ya había producido, por lo que \emph{no necesita SAM,
ni el checkpoint de 375~MB, ni GPU}.

\begin{longtable}{@{}p{4.2cm} p{11.4cm}@{}}
\toprule
\textbf{Módulo} & \textbf{Contenido} \\
\midrule
\endhead
\gh{Geometry/lorentz.py} & Álgebra de Minkowski, marcos $\eta$-ortonormales, el boost que
fija la esfera, acción proyectiva y detección del plano evanescente. \\
\gh{Geometry/similarity.py} & Grupo de similaridades, el par (esfera, eje) como estructura
de datos, el invariante \Rd{} y el contraejemplo. \\
\gh{Geometry/sphere\_axis.py} & Extracción del par desde malla y desde CSV; estimador
robusto del eje; incertidumbre por bootstrap; calibración del sesgo por truncamiento. \\
\gh{Geometry/stl\_io.py} & Lectura y escritura de STL binario y ASCII con
\texttt{numpy}; muestreo de superficie ponderado por área. \\
\gh{Geometry/transfer.py} & Forma canónica, composición completa, transporte de mallas y
métricas de eficacia. \\
\gh{Geometry/selftest.py} & Las 16 pruebas de verificación. \\
\gh{exercises\_geometry.py} & Interfaz de línea de comandos. \\
\bottomrule
\end{longtable}

\subsection{La composición}

La transformación se construye por conjugación:
\[
\Phi \;=\; K_{\text{pac}}^{-1} \circ M \circ K_{\text{std}},
\]
donde $K$ lleva un par a la forma canónica ---esfera de radio $\rho$ en el origen, eje
paralelo a $e_1$ a altura $h = \rho\, d/R$--- y $M$ es el boost de \Ot{} que fija esa esfera
y manda la cuerda canónica del referente en la del paciente. Con
$Q = [v\;\; n\;\; w]$ por columnas ($v$ la dirección del eje, $w$ el unitario del centro
hacia el pie de la perpendicular, $n = w \times v$), se toma $K(x) = (\rho/R)\,Q^{\mathsf T}(x - C)$.

Todo se representa en coordenadas homogéneas en la carta $w = \rho$, de modo que $\Phi$ es
\textbf{una sola matriz $4\times4$}: una única deshomogeneización y sin pérdida de precisión
por composiciones sucesivas.

Si las razones \Rd{} coinciden, entonces $h_{\text{std}} = h_{\text{pac}}$, sale $M = I$ y
$\Phi$ se reduce a una similaridad. En caso contrario $\Phi$ es genuinamente proyectiva:
preserva rectas, planos y razones dobles, pero \emph{no} ángulos ni razones de distancias.
Como una proyectividad manda planos en planos, los triángulos siguen siendo triángulos
planos y las mallas se transportan vértice a vértice sin remallar.

% ===========================================================================
\section{Verificación}

El núcleo se verifica contra el resultado publicado en el enunciado. Con $\rho = 5$,
$A_1 = (-\sqrt{21},0,2)$, $B_1 = (\sqrt{21},0,2)$, $A_2 = (-4,0,-3)$, $B_2 = (4,0,-3)$, la
implementación reproduce

\[
M = \begin{pmatrix}
1 & 0 & 0 & 0\\
0 & 1 & 0 & 0\\
0 & 0 & \tfrac{31\sqrt{21}}{84} & -\tfrac{25\sqrt{21}}{84}\\[2pt]
0 & 0 & -\tfrac{25\sqrt{21}}{84} & \tfrac{31\sqrt{21}}{84}
\end{pmatrix}
\]

entrada por entrada, con $\lVert M^{\mathsf T}\eta M - \eta\rVert = 4.4\times10^{-16}$ y factor
proyectivo $4/\sqrt{21} = 0.872871560944$ exacto.

\begin{center}
\begin{tabular}{@{}p{7.6cm} p{7.6cm}@{}}
\toprule
\textbf{Prueba} & \textbf{Resultado} \\
\midrule
$O(3) \subset \Ot$, 200 rotaciones & $\lVert M^{\mathsf T}\eta M - \eta\rVert \le 7.8\times10^{-16}$ \\
Matriz publicada, entrada por entrada & $4.4\times10^{-16}$ \\
La esfera se preserva, 500 puntos & $8.9\times10^{-15}$ mm \\
La cuerda va a la cuerda & colinealidad $1.3\times10^{-14}$ \\
Generaliza a 30 cuerdas aleatorias & $\le 2.1\times10^{-12}$ \\
Coincide con la receta del enunciado & $\le 1.1\times10^{-13}$, y $M_0 = M$ \\
Normalización \emph{timelike} sin \texttt{NaN} & $\langle c,c\rangle = -21 < 0$ \\
Familia de gauge & cortes fijos, deslizamiento interior \\
Similaridad de la diapositiva 12 & $\det R = +1$, $\lVert R^{\mathsf T}R - I\rVert = 3.9\times10^{-15}$ \\
Invariante \Rd & hueco nulo cuando las razones coinciden \\
Eje paralelo, 50 pares aleatorios & paralelismo $\le 1.2\times10^{-15}$ \\
Contraejemplo & $5/3 \ne 5$, hueco irreducible 2.0 mm \\
\bottomrule
\end{tabular}
\end{center}

\noindent Total: \textbf{16 de 16}. Se ejecutan con
\texttt{python exercises\_geometry.py -{}-selftest}.

\subsection{Dos rectificaciones}

\paragraph{La receta del enunciado es correcta.} Una primera lectura de este trabajo
concluyó que la receta de las diapositivas 20--21 no generalizaba. \textbf{Era un error de
transcripción}: los vectores $s_{20}, t_{20}$ pertenecen al segundo marco, no al primero.
Con la lectura correcta la receta se mantiene en \Ot{} sobre cuerdas arbitrarias y produce
exactamente la misma matriz que la ruta corta $M = F_2\,(\eta F_1^{\mathsf T} \eta)$.

Lo que sí queda como aporte: los cuatro vectores escritos a mano en la diapositiva son la
completación $\eta$-ortonormal de cada marco (hay que \emph{generarlos} en el caso general);
el paso $s_2 = M_0 s_1$ es redundante, pues sale $M_0 = M$; y el vector
$c = (n_A + n_B)/2$ es \emph{timelike}, con norma de Minkowski negativa, de modo que
normalizarlo sin cuidado produce \texttt{NaN}.

\paragraph{Un fallo silencioso que conviene conocer.} El error cometido produce una matriz
que \emph{sigue mandando $A_1 \to A_2$ y $B_1 \to B_2$ con error $4.4\times10^{-16}$ y aun
así está fuera de \Ot} ($\lVert M^{\mathsf T}\eta M - \eta\rVert = 1.88$). Mapea bien y está
mal. Por eso la implementación verifica la pertenencia al grupo antes de devolver nada, y
hay una prueba dedicada a ese caso.

% ===========================================================================
\section{Resultados}

\subsection{El par del húmero de referencia}

De \texttt{Humero.stl} (2590 triángulos, 155 mm de longitud):
$R = 22.89$ mm, $d = 10.50$ mm, $\Rd = 2.18$ ($d/R = 0.459$), con RMS del eje de 0.75 mm.
El enunciado anticipa $R \approx 24$ mm y $d/R \approx 0.4$: el referente cae donde debe.

Tres precauciones resultaron imprescindibles para llegar ahí, y todas son del estimador,
no del dato:

\begin{itemize}[leftmargin=1.2em,itemsep=1pt]
\item \textbf{Muestrear la superficie, no los vértices.} El 70\,\% de los 1301 vértices están
en la cabeza y la diáfisis se queda con unos 10 por banda. Muestreando 40\,000 puntos
ponderados por área, el RMS del eje baja de 2.09 mm a 0.75 mm.
\item \textbf{Rechazar losas por radio anómalo.} La malla tiene geometría espuria en algunas
bandas (radios transversales de 32 a 56 mm donde la diáfisis mide 11). Sin ese filtro el eje
salía con $d = 49.5$ mm.
\item \textbf{No reseleccionar la región diafisaria iterando} sobre el eje ya refinado: se
realimenta y diverge ($8.5\degree \to 8.4\degree \to 11.1\degree$).
\end{itemize}

El eje correcto queda a $6.1\degree$ del PCA global, que es exactamente lo que la cabeza tuerce.

\subsection{Los pares de los pacientes}

De los 8 datasets, 5 tienen los CSV necesarios y \textbf{3 superan los criterios de calidad},
fijados a priori en el código y reportados como banderas, no aplicados a ojo.

\begin{center}
\begin{tabular}{@{}l r r r r r c@{}}
\toprule
\textbf{Dataset} & \textbf{$R$ (mm)} & \textbf{$d$ (mm)} & \textbf{\Rd} &
\textbf{$\sigma$ eje} & \textbf{art./diáf.} & \textbf{Usable} \\
\midrule
\texttt{pd\_tse\_fs\_tra\_320\_fov150\_4}  & 24.56 & 4.77 & 5.15 & $1.02\degree$ & 6/11 & sí \\
\texttt{t1\_tse\_fs\_tra\_320\_fov150\_11} & 23.92 & 3.26 & 7.33 & $1.21\degree$ & 7/10 & sí \\
\texttt{t1\_tse\_fs\_tra\_320\_fov150\_5}  & 23.95 & 4.43 & 5.41 & $0.36\degree$ & 8/9  & sí \\
\texttt{AXIAL-P5SE1 \ldots SIEMENS}        & 24.57 & 3.16 & 7.78 & $3.13\degree$ & 3/11 & no \\
\texttt{AXIAL fs PD FSE \ldots GE}         & 19.69 & 1.54 & 12.78 & $2.30\degree$ & 6/10 & no \\
\bottomrule
\end{tabular}
\end{center}

\noindent\small
Motivo de exclusión: el dataset SIEMENS tiene solo 3 círculos articulares y su eje no
converge; el dataset GE da un radio de 19.69 mm, fuera del rango anatómico admitido.
\normalsize

\subsection{El hallazgo principal: la violación de \texorpdfstring{\Rd}{R/d} es un artefacto}

La lectura inmediata de la tabla anterior es que el referente tiene $\Rd = 2.18$ y los
pacientes entre 5.15 y 12.78, luego la condición de la diapositiva 13 se viola por un factor
de 2.5 a 6 y \Ot{} es imprescindible. \textbf{Esa conclusión no se sostiene.}

Los volúmenes son cortos (unos 68 mm, de 19 a 31 cortes) y solo cubren el húmero proximal:
los cortes que el pipeline toma por diáfisis son en realidad cuello quirúrgico y metáfisis.
El propio enunciado ya lo anotaba como pendiente en la diapositiva 16.

Como el húmero de referencia \emph{sí} está completo, el sesgo es medible: se trunca su malla
a la misma extensión axial y se re-estima con el \emph{mismo} estimador.

\begin{center}
\begin{tabular}{@{}r r r r r@{}}
\toprule
\textbf{Extensión} & \textbf{Error angular} & \textbf{$d$ leído} &
\textbf{$d_{\text{esc}}$} & \textbf{\Rd{} leído} \\
\midrule
55 mm  & $55.5\degree$ & 1.71 mm  & 0.163 & 12.43 \\
60 mm  & $19.1\degree$ & 4.50 mm  & 0.428 & 5.02 \\
\textbf{68 mm} & $\mathbf{10.0\degree}$ & \textbf{4.78 mm} & \textbf{0.455} & \textbf{4.80} \\
80 mm  & $10.0\degree$ & 3.70 mm  & 0.352 & 6.24 \\
100 mm & $6.3\degree$  & 4.87 mm  & 0.464 & 4.66 \\
130 mm & $3.4\degree$  & 6.49 mm  & 0.618 & 3.50 \\
\midrule
\textbf{completa} & --- & \textbf{10.50 mm} & 1.000 & \textbf{2.18} \\
\bottomrule
\end{tabular}
\end{center}

Es decir: \textbf{el propio húmero de referencia, visto con la cobertura de una resonancia,
aparenta $\Rd \approx 5.2$ cuando en realidad vale 2.18.} Aplicando el factor
$d_{\text{esc}} = 0.425 \pm 0.044$ (ventana de 58 a 100 mm):

\begin{center}
\begin{tabular}{@{}l r r r@{}}
\toprule
\textbf{Dataset} & \textbf{\Rd{} leído} & \textbf{\Rd{} corregido} & \textbf{Referente} \\
\midrule
\texttt{pd\_tse\_fs\_tra\_320\_fov150\_4}  & 5.15 & \textbf{2.19} & 2.18 \\
\texttt{t1\_tse\_fs\_tra\_320\_fov150\_5}  & 5.41 & \textbf{2.30} & 2.18 \\
\texttt{t1\_tse\_fs\_tra\_320\_fov150\_11} & 7.33 & \textbf{3.11} & 2.18 \\
\bottomrule
\end{tabular}
\end{center}

Dos de los tres datasets caen prácticamente sobre el referente. \textbf{Con estos datos no
se puede afirmar que haga falta \Ot.} La maquinaria hay que construirla igual ---el enunciado
la pide y el caso general la necesita--- pero la conclusión honesta es que la violación
observada era, en su mayor parte, el sesgo del eje.

\paragraph{Precisión no es exactitud.} El bootstrap da $\sigma$ del eje de apenas
$0.36\degree$ a $1.21\degree$: el eje está muy bien \emph{determinado} por los datos
disponibles. Pero el truncamiento introduce $10\degree$ de error \emph{sistemático}. Reportar
solo $\sigma$ habría dado una falsa sensación de exactitud.

\subsection{Transporte y validación cruzada}

$\Phi$ verifica numéricamente en los 3 datasets: la esfera del referente va sobre la del
paciente con error máximo de $7.8\times10^{-14}$ mm, el eje sobre el eje con
$5.6\times10^{-13}$ mm, y $M$ está en \Ot{} a $10^{-16}$. Las cuatro mallas (húmero,
infraespinoso, redondo menor y subescapular; entre 2590 y 38\,776 triángulos) se transportan
sin cruzar el plano evanescente.

La comparación entre usar el $d$ crudo y el corregido es la validación cruzada del hallazgo:
la corrección se deriva del \emph{referente} y se valida contra la nube de borde del
\emph{paciente}, que no intervino en derivarla.

\begin{center}
\begin{tabular}{@{}l c r r r r@{}}
\toprule
\textbf{Dataset} & \textbf{Variante} & \textbf{\Rd} &
\textbf{$w_{\max}/w_{\min}$} & \textbf{RMS} & \textbf{P95} \\
\midrule
\multirow{2}{*}{\texttt{pd\_\ldots\_fov150\_4}}
 & crudo     & 5.15 & 3.39 & 4.27 mm & 8.76 mm \\
 & corregido & 2.19 & \textbf{1.01} & \textbf{3.33 mm} & \textbf{6.84 mm} \\
\addlinespace
\multirow{2}{*}{\texttt{t1\_\ldots\_fov150\_11}}
 & crudo     & 7.33 & 5.00 & 4.47 mm & 9.61 mm \\
 & corregido & 3.11 & \textbf{1.77} & \textbf{3.60 mm} & \textbf{7.60 mm} \\
\addlinespace
\multirow{2}{*}{\texttt{t1\_\ldots\_fov150\_5}}
 & crudo     & 5.41 & 3.58 & 4.40 mm & 9.06 mm \\
 & corregido & 2.30 & \textbf{1.10} & \textbf{3.07 mm} & \textbf{6.18 mm} \\
\bottomrule
\end{tabular}
\end{center}

La columna $w_{\max}/w_{\min}$ mide la distorsión proyectiva: cae de 3.4--5.0 a 1.0--1.8, es
decir el transporte pasa de ser fuertemente proyectivo a ser casi una similaridad. Y la
fidelidad del hueso mejora en los tres casos, de forma consistente y sin haber ajustado nada
contra esos datos.

\subsection{Eficacia de modelación}

\begin{center}
\colorbox{grisclaro}{\parbox{0.94\textwidth}{\small
\textcolor{rojoalerta}{\textbf{Advertencia.}} No existe segmentación muscular del paciente.
Ninguna de estas métricas mide exactitud anatómica de los músculos: miden
\emph{consistencia geométrica} del transporte. Índices de Dice o IoU musculares, espesor
tendinoso y detección de desgarro son sencillamente no evaluables con los datos
disponibles, y no se sustituyen por nada.
}}
\end{center}

\paragraph{M1, fidelidad ósea.} Es la única métrica con verdad de referencia, porque el
húmero del referente sí tiene contraparte en el paciente. Distancia de la nube de borde de
la resonancia a la superficie del húmero transportado, restringida a la banda que la
resonancia realmente cubre: RMS de 3.07 a 3.60 mm, mediana de 2.13 a 2.33 mm (variante
corregida). Como los músculos viajan solidarios al hueso por la misma $\Phi$, este error es
una \emph{cota inferior} del error muscular.

\paragraph{M2, plausibilidad anatómica.} La penetración de los músculos en el hueso es
inferior al 0.1\,\% de los vértices para el infraespinoso y el redondo menor, y del 6.6\,\%
para el subescapular. Las huellas de inserción se sitúan de forma muy estable entre datasets
($\theta \approx 45$--$58\degree$ para el infraespinoso, $74$--$82\degree$ para el redondo
menor, $207\degree$ para el subescapular), lo que indica que el transporte es reproducible.

\paragraph{Sobre los sectores angulares.} Estas cifras \emph{no} deben compararse todavía con
la regla nominal de la presentación. El origen angular $\theta = 0$ del pipeline procede de
una base obtenida por SVD, sin anclaje anatómico, de modo que los sectores nominales no están
bien definidos hasta fijar una dirección de referencia. Que el trabajo lo haga evidente ya es
un resultado: la regla de sectores necesita un origen anatómico explícito antes de poder
validarse.

% ===========================================================================
\section{Limitaciones}

\begin{itemize}[leftmargin=1.2em,itemsep=2pt]
\item \textbf{Cobertura axial insuficiente.} Es la limitación dominante y la causa del sesgo
del eje. La solución real es cobertura distal: o un stack axial que llegue a la mitad del
brazo, o fusionar la esfera del axial con el eje de un stack \emph{coronal}, donde la
diáfisis se ve a lo largo dentro de cada imagen. Existe uno sin procesar en el repositorio.
\item \textbf{Tres datasets válidos.} No es una validación, es una prueba de concepto. No se
promedian entre sí.
\item \textbf{Falta el supraespinoso}, justo el tendón donde ocurre la mayoría de los
desgarros y el sector central de la hipótesis de la presentación. No se fabrica ni se estima
por complemento.
\item \textbf{La corrección del sesgo es de primer orden}: corrige la magnitud de $d$, no el
error angular de $10\degree$, y el factor tiene una dispersión de en torno al 10\,\%.
\item \textbf{La libertad de gauge no se ha explotado.} El boost admite una familia de un
parámetro que fija esfera y eje pero desliza los puntos a lo largo de la recta; está
implementada y verificada, pero optimizarla contra la nube del paciente queda pendiente.
\end{itemize}

% ===========================================================================
\section{Cómo reproducir}

\begin{verbatim}
git clone https://github.com/gustavop-dev/humerous-detection-fat-supress.git
cd humerous-detection-fat-supress
git checkout geometry-exercises

python3 -m venv venv-analysis
source venv-analysis/bin/activate
pip install numpy scipy matplotlib     # sin torch, sin SAM, sin GPU

python exercises_geometry.py --selftest    # 16/16
python exercises_geometry.py --all         # pipeline completo
\end{verbatim}

\noindent
Salidas por dataset en \texttt{Datasets/Out/<dataset>/geometry/}:
\texttt{patient\_frame.json}, \texttt{transfer\_report.json} y
\texttt{muscles\_mapped/*.stl}. Salidas globales en \texttt{excercises/out/}:
\texttt{standard\_frame.json}, \texttt{axis\_bias\_calibration.json} y
\texttt{bias\_comparison.json}.

\vspace{0.4cm}
\hrule
\vspace{0.25cm}
\noindent\small
Código y documentación: \url{\ghrepo} \quad rama \texttt{\ghbranch}.

\end{document}
